\documentclass[10pt,a4paper]{article}
\usepackage[margin=2cm]{geometry}
\usepackage{amsmath,amssymb,amsthm,graphicx,color,xcolor,xspace,stmaryrd}
\usepackage[bookmarks,bookmarksopen]{hyperref}
\usepackage{paralist,tikz,todonotes}
\usetikzlibrary{arrows.meta,positioning}
\tikzset{ew/.style={to path={(\tikztostart.east)--(\tikztotarget.west)}}}
\newtheorem{theorem}{Theorem}
\newtheorem{lemma}[theorem]{Lemma}
\newtheorem{proposition}[theorem]{Proposition}
\newtheorem{corollary}[theorem]{Corollary}
\newtheorem{definition}[theorem]{Definition}
\newtheorem{example}[theorem]{Example}
\newtheorem{remark}[theorem]{Remark}
\AtBeginDocument{\pagestyle{plain}}
\makeatletter
% highlighting of changes
% new in mono DHOL over HOL
\newcommand{\hlm}[1]{{\color{blue}#1}}
% new in poly DHOL
\newcommand{\hlp}[1]{{\color{red}#1}}
% new with dependent type variables
\newcommand{\hldtv}[1]{{\color{purple}#1}}

% syntax

\newcommand{\type}{\mathtt{type}}
\newcommand{\bool}{o}
\newcommand{\impl}{\Rightarrow}

% term equality at a type
\newcommand{\eq}[3]{#2=_{#1}#3}
% typeness
\newcommand{\tp}[1]{#1\,:\,\type}
% type equality
\newcommand{\teq}[2]{#1\equiv#2}

% Pi with a list of bindings
\newcommand{\Pg}[1]{\Pi_{#1}}
% binding given as name:type
\renewcommand{\P}[2]{\Pg{#1:#2}}
% binding of type variables
\newcommand{\Pt}[1]{\Pg{#1}}
% lambda with a list of bindings
\newcommand{\lamg}[1]{\lambda_{#1}}
% lambda binding name:type
\newcommand{\lam}[2]{\lamg{#1:#2}}

% keyword to introduce an assumption in theory or context
\newcommand{\ass}{\triangleright\;}

% sequences
\newcommand{\V}[1]{\vec{#1}}
\newcommand{\vA}{\V A}
\newcommand{\valpha}{\V \alpha}
\newcommand{\valphaper}{\V {\alpha^*}}
\newcommand{\vt}{\V t}
\newcommand{\vx}{\V x}

% judgments
\newcommand{\Thy}[1]{\vdash #1\;\mathtt{Thy}}
% all judgments below use T as default value for the theory
\newcommand{\Ctx}[2][T]{\vdash_{#1} #2\;\mathtt{Ctx}}
\newcommand{\Sub}[4][T]{#2\vdash_{#1} #4\;\leftarrow\; #3}
\newcommand{\Type}[3][T]{#2\vdash_{#1} \tp{#3}}
\newcommand{\Term}[4][T]{#2\vdash_{#1} #3\,:\,#4}
\newcommand{\TEq}[4][T]{#2\vdash_{#1} \teq{#3}{#4}}
\newcommand{\SubEq}[5][T]{#2\vdash_{#1} #4\equiv_{#3}#5}
\newcommand{\True}[3][T]{#2\vdash_{#1} #3}

% abbreviations for making \Gamma the default context
\newcommand{\SubG}[3][\Gamma]{\Sub{#1}{#2}{#3}}
\newcommand{\SubEqG}[4][\Gamma]{\SubEq{#1}{#2}{#3}{#4}}
\newcommand{\TypeG}[2][\Gamma]{\Type{#1}{#2}}
\newcommand{\TermG}[3][\Gamma]{\Term{#1}{#2}{#3}}
\newcommand{\TEqG}[3][\Gamma]{\TEq{#1}{#2}{#3}}
\newcommand{\TrueG}[2][\Gamma]{\True{#1}{#2}}

% rules

% \rul[macroname]{rulename}{hyps}{conc}; defines \macroname to produce "Rule rulename"
\newcommand{\rul}[4][]{\hypertarget{#1}{\cfrac{#3}{#4}{\scriptstyle#2}}\ifnonempty{#1}{\expandafter\gdef\csname#1\endcsname{\hyperlink{#1}{Rule #2}\xspace}}}


% \ifnonempty{a}{b} produces b if a is non-empty
\newcommand{\ifnonempty}[3][]{\def\@empty{}\def\@test{#2}\ifx\@test\@empty#1\else#3\fi}

% translation

% dependency erasure
\newcommand{\tr}[1]{\overline{#1}}
%  PER of a type
\newcommand{\per}[1]{{#1}^\ast}
% \perapp Atu: the PER for A applied to t and u
\newcommand{\perapp}[3]{\per {#1}\ #2\ #3}
% formula to say that a relation is PER
\newcommand{\isper}[1]{\mathrm{PER}(#1)}

% only type variables
\newcommand{\ot}[1]{#1^{\textcolor{blue}{y}}}
% only type and term variables
\newcommand{\ott}[1]{#1^{\textcolor{red}{ye}}}

% judgments cont.
\newcommand{\CtxT}[2][\tr{T}]{\vdash_{#1} #2\;\mathtt{Ctx}}
\newcommand{\SubT}[4][\tr{T}]{#2\vdash_{#1} #4\leftarrow #3}
\newcommand{\TypeT}[3][\tr{T}]{#2\vdash_{#1} \tp{#3}}
\newcommand{\TermT}[4][\tr{T}]{#2\vdash_{#1} #3:#4}
\newcommand{\TEqT}[4][\tr{T}]{#2\vdash_{#1} \teq{#3}{#4}}
\newcommand{\SubEqT}[5][\tr{T}]{#2\vdash_{#1} \teq{#4}{#5}}
\newcommand{\TrueT}[3][\tr{T}]{#2\vdash_{#1} #3}

%examples
\newcommand{\cn}[1]{\mathtt{#1}}
\newcommand{\nat}{\cn{nat}}
\newcommand{\nil}{\cn{nil}}
\newcommand{\suc}{\cn{suc}}
\newcommand{\+}{\cn{+}}
\newcommand{\vect}{\cn{vec}}
\newcommand{\cons}{\cn{cons}}
\newcommand{\app}{\cn{++}}
\newcommand{\fin}[1]{\cn{fin}\ #1}
\newcommand{\fnext}{\cn{fnext}}
\newcommand{\ftop}{\cn{ftop}}
\newcommand{\Hlist}[2]{\cn{Hlist}\ #1\ #2}
\newcommand{\hlist}{\cn{hlist}}
\newcommand{\hget}{\cn{hget}}
\newcommand{\clr}{\cn{color}}
\newcommand{\bl}{\cn{black}}
\newcommand{\rd}{\cn{red}}
\newcommand{\rt}{\cn{red\_tree}}
\newcommand{\bt}{\cn{black\_tree}}
\newcommand{\tree}{\cn{tree}}
\newcommand{\leaf}{\cn{leaf}}
\newcommand{\rev}{\cn{rev}}

\newcommand{\ha}{\hlp{\alpha}}
\newcommand{\hn}{\hlm{n}}
\newcommand{\hm}{\hlm{m}}

% better figure placement
\renewcommand{\topfraction}{0.95}	% max fraction of floats at top
\setcounter{topnumber}{2}
\setcounter{bottomnumber}{2}
\setcounter{totalnumber}{4}
\renewcommand{\textfraction}{0.07}	% allow minimal text w. figs
\renewcommand{\floatpagefraction}{0.9}	% require fuller float pages

% makes strings into macros for themselves so that we can control the font
\newcommand{\makecn}[1]{\expandafter\def\csname#1\endcsname{\mathtt{#1}}}
% allows writing function application as f"{x,y,z}
\newcommand{\fold}[2]{\def\@tmpop{\relax}\@for\@I:=#2\do{\@tmpop\@I\def\@tmpop{#1}}}
\catcode`"=\active
\def"#1{\,\fold{\,}{#1}}

\makeatother
